% Transcription — fonds Alexandre Grothendieck, shelfmark 162-1, batch 1 (pages 1-9).
% Established from the facsimile archives/batches/162-1/batch-01.pdf (a mirror of
% https://grothendieck.umontpellier.fr/162-1.pdf), per the skill
% transcribe-grothendieck.
% The folder is a lending register, not a mathematical manuscript: a borrower's
% name at the left, what they took at the right. Most entries are scored
% through, which on a register of this kind reads as a return rather than as a
% change of mind — but the page nowhere says so, so the strokes are transcribed
% as strokes and the reading is left to stand.
% A large part of the register is illegible: the strikings are heavy, several
% run over one another, and a name struck twice is often gone. Everything
% doubtful is marked and nothing has been supplied by conjecture. The proportion
% of \ill here is far higher than in a mathematical folder, and that is the
% honest state of the document.
% Pass: Opus 5 (claude-opus-5), 2026-08-16 — first pass, unchecked against the pages by a human.
\documentclass[11pt,a4paper]{article}
% Le préambule commun à toutes les transcriptions du fonds Grothendieck.
%
% Il tient en une idée : **l'incertitude se compose, elle ne se résout pas.**
% Une transcription qui lisse un mot douteux a détruit la seule chose pour
% laquelle on la faisait — un lecteur doit pouvoir savoir, à chaque endroit,
% ce qui était sur le feuillet et ce qui a été ajouté. Les six macros qui
% suivent sont donc l'appareil critique, et non de la décoration.
%
% Les mêmes commandes sont reconnues par `scripts/render.mjs`, qui produit la
% vue de lecture du site. Une macro ajoutée ici doit l'être aussi là-bas,
% faute de quoi le rendu échouera bruyamment — ce qui est voulu : un rendu
% silencieusement faux est pire qu'un rendu absent.

\ProvidesPackage{grothendieck}[2026/08/08 Transcriptions du fonds Grothendieck]

\usepackage{amsmath,amssymb,amsthm}
\usepackage{mathrsfs}
\usepackage{stmaryrd}
% Les diagrammes commutatifs sont partout dans ce fonds : c'est la langue
% même de la géométrie algébrique de Grothendieck, et une transcription qui
% les rendrait en prose ne transcrirait rien.
\usepackage{tikz-cd}
% Français par défaut — c'est la langue des feuillets — mais l'anglais est
% chargé aussi : la traduction et le résumé sont des documents anglais, et la
% typographie française (espaces avant les deux-points) y serait une faute.
% Ces documents basculent par \selectlanguage{english} en tête de corps.
\usepackage[english,french]{babel}
\usepackage{fontspec}
\usepackage{geometry}
\geometry{a4paper,margin=25mm,marginparwidth=28mm}
\usepackage{marginnote}
\usepackage{xcolor}
% ulem, not soul. `soul`'s \st and \ul are built on a character-by-character
% scan that XeLaTeX's font machinery breaks: the document fails with an
% undefined control sequence rather than degrading. ulem does the same two jobs
% and is Unicode-engine safe. `normalem` stops it hijacking \emph, which the
% transcriptions use for Grothendieck's own emphasis.
\usepackage[normalem]{ulem}
\usepackage{hyperref}
\hypersetup{colorlinks=true,linkcolor=black,urlcolor=black,pdfborder={0 0 0}}

% --- Métadonnées du lot -----------------------------------------------------
% Reprises telles quelles par le script de rendu, qui les affiche en tête de
% la vue de lecture. Elles fixent ce que le fichier prétend couvrir : sans
% elles, une transcription flotte sans point d'attache dans un fonds de seize
% mille feuillets.

\newcommand{\folder}[1]{\gdef\grothFolder{#1}}
\newcommand{\batch}[1]{\gdef\grothBatch{#1}}
\newcommand{\leaves}[2]{\gdef\grothFirst{#1}\gdef\grothLast{#2}}
\newcommand{\foldertitle}[1]{\gdef\grothFoldertitle{#1}}
\gdef\grothFolder{?}\gdef\grothBatch{?}\gdef\grothFirst{?}\gdef\grothLast{?}\gdef\grothFoldertitle{}

% --- Le résumé --------------------------------------------------------------
%
% Il ouvre la lecture modernisée, sous le titre « Résumé », et il est composé
% en petit. Ce n'est pas un ornement : c'est le seul passage du document qui
% ne soit pas des mathématiques, et un lecteur doit pouvoir le reconnaître
% avant de l'avoir lu — puis le sauter s'il connaît déjà le sujet, ou s'y
% arrêter s'il ne le connaît pas. Le filet qui le suit marque la reprise du
% corps.

\newenvironment{resume}
  {\small\setlength{\parskip}{0.45em}}
  {\par\medskip\hrule\medskip}

% --- Mots-clés ---------------------------------------------------------------
%
% En anglais, en clôture du résumé de la lecture modernisée : le vocabulaire
% moderne sous lequel chercher ce que les feuillets construisent. Le manifeste
% du site (`npm run manifest`) les extrait pour étiqueter la cote entière —
% cette ligne est la source unique des tags, il n'y a pas de fichier à côté.
\newcommand{\keywords}[1]{\par\smallskip\noindent
  {\footnotesize\textsc{Keywords} --- \textit{#1}}\par}

% --- L'appareil critique ----------------------------------------------------

% Le numéro de feuillet, en marge. C'est l'ancre qui permet de régler un
% désaccord sur une formule en regardant *un* feuillet plutôt qu'une chemise
% de six cents. Le site s'en sert aussi pour tourner les pages du fac-similé
% quand on fait défiler la transcription.
\newcommand{\leaf}[1]{%
  \marginnote{\footnotesize\color{gray!70}\textsf{#1}}%
  \label{leaf:#1}\ignorespaces}

% Illisible : on ne devine pas. Un mot inventé qui se lit comme les autres est
% le pire résultat possible.
\newcommand{\ill}{\textcolor{red!60!black}{[\,\ldots\,]}}

% Lecture douteuse : le mot est donné, et signalé comme incertain.
\newcommand{\uncertain}[1]{\uline{#1}}

% Ajout de l'éditeur — une lettre manquante, un mot sous-entendu.
\newcommand{\add}[1]{\textcolor{blue!50!black}{[#1]}}

% Biffé par Grothendieck lui-même. Ce qu'il a rayé fait partie du document :
% une rature dit souvent où la pensée a changé de direction.
\newcommand{\struck}[1]{\sout{#1}}

% Note du transcripteur, sur l'état matériel du feuillet.
\newcommand{\note}[1]{\par\smallskip{\small\color{gray!60!black}\textit{[#1]}}\par\smallskip}

% Note marginale de Grothendieck, distincte du corps du texte.
\newcommand{\marginal}[1]{\par\smallskip\noindent\hspace{1em}%
  {\small\color{gray!60!black}#1}\par\smallskip}

% --- Titre ------------------------------------------------------------------

\AtBeginDocument{%
  \begin{center}
    {\large\bfseries Fonds Alexandre Grothendieck --- cote n\textsuperscript{o}~\grothFolder}\\[2pt]
    {\itshape\grothFoldertitle}\\[4pt]
    {\small feuillets \grothFirst--\grothLast\ (lot \grothBatch)}\\[2pt]
    {\footnotesize\color{gray!60!black}Transcription établie d'après le fac-similé numérisé
     par l'Université de Montpellier.\\
     Les lectures incertaines sont soulignées, les illisibles notées [\,\ldots\,],
     les ajouts entre crochets.}
  \end{center}
  \vspace{1em}}

\endinput


\folder{162-1}
\batch{1}
\pages{1}{9}
\dating{[à partir de 1967] — le groupe « Dossiers rassemblés par Grothendieck sur différentes thématiques » (161-1 à 162-6) est daté [après 1961-vers 1977]}
\watermark{Édition de démonstration}
\foldertitle{Papiers prêtés [liste de noms et de références documentaires] : notes manuscrites (s.d.)}

\begin{document}

\page{1}
Titre porté à l'encre en haut à droite de la chemise : « Papiers prêtés ». Au
centre, au crayon, un chiffre romain \uncertain{III} et, au-dessous, une
indication entre parenthèses restée illisible : \ill

\note{le registre qui suit se lit en deux colonnes — l'emprunteur à gauche,
souvent souligné, et parfois suivi d'un chiffre ; ce qu'il a emporté à droite.
Ce chiffre, 1 ou 2, revient trop régulièrement pour être un nombre de pièces ;
il compte vraisemblablement les rappels, mais le registre ne l'explique pas et
on ne le suppose pas ici}

\page{2}

\textbf{Deligne} \emph{(suite)} — \struck{Hironaka, thèse} ;
\struck{\uncertain{Griffiths} périodes} ; \struck{SGA 6 XIV} ; \struck{Chern et
classes de Chern des représentations} ; \struck{Quillen sur \uncertain{sém.} de
Adams} ; \struck{2 lettres à \ill\ sur Atiyah--Adams)} \add{du} 10.I.66 ;
\struck{thèse \ill} ; Quillen, \emph{Derivations} \ill ; \struck{Hartshorne \ill\
cohomologie \ill}

\note{le « 10.I.66 » est la seule date portée sur cette page ; elle est
antérieure au « [à partir de 1967] » de l'inventaire, qui a dû être déduit
d'ailleurs — la lettre du 21.9.67 portée p.~5 en est le candidat le plus
proche}

\textbf{\struck{Giraud}} — \struck{Notes sur \ill\ 67}

\textbf{\struck{\ill}} — \struck{Lettre sur Hartshorne} ; \struck{Notes Algèbre
III Bourbaki}

\textbf{Shih} \emph{1} — R. Wood à Bourbaki, calcul différentiel ;
\struck{3 articles Schwarzenberger sur fibrés vectoriels} ; Résumé PTT Grenoble

\textbf{\struck{\ill}} — \struck{\ill} Papiers Atiyah, topologie et opérateurs
elliptiques

\textbf{\struck{Bertin}} — \struck{EGA IV 16--19, photocopie}

\textbf{\struck{Mme Hakim}} — \struck{Séminaire Gabriel, deux séminaires,
exp.~7}

\textbf{Ferrand} \emph{1} — \struck{EGA IV 20, 21} \struck{\ill} ; Lettre à
Dieudonné sur \ill\ de dimension 2 ; \struck{Notes sur R.R.} \add{entouré, suivi d'un point d'interrogation} ; Notes sur \ill\ dét.\ / sur cat.\ de Picard
\add{et} Godement ; \struck{Papiers \uncertain{Hertzbrouck}}

\textbf{Jouanolou} \emph{1} — Papiers Hirzebruch (congrès d'Amsterdam ?) ;
Lettre : Serre \add{sur} \ill\ fonctionnelles, suite, début de monodromie \ill

\textbf{\struck{Berthelot}} — \struck{Notes sur stratifications} ; \struck{lettre
Deligne sur opérateurs elliptiques} ; \struck{Notes R.R.\ (copie de Serre)} ;
\struck{Notes Serre sur \ill} ; \struck{Calcul \ill} ; \struck{Rapport Serre,
représentations de groupes finis}

\textbf{\struck{Kleiman}} — \struck{Manuscrits ? EGA IV 20} ; \struck{EGA V} ;
\struck{Notes Chevalley sur groupes \ill} ; \struck{Notes \ill} ;
\emph{Problèmes de Géométrie Algébrique}

\textbf{\struck{Verdier}} — SGA 4 fasc.~1, \ill\ ; \struck{Gabriel}

\textbf{\struck{\ill}} — \struck{Article \uncertain{Grünberg} sur \ill}

\page{3}

\textbf{Raynaud} — Mumford, \emph{$\Theta$ foundations} \add{I}, II ; Papiers
Kiehl ; Artin--Winters

\textbf{Illusie} \emph{1} — Monsky--Washnitzer I, II ; Lettres Hartshorne ;
Hartshorne, \emph{Affine duality} ; Lettre Quillen ; Notes sur extensions
infinitésimales ; \struck{Notes sur tapis (Quillen)} ; R. Wood à Bourbaki,
calcul différentiel ; Complexe cotangent relatif de Quillen ; Notes sur
l'extension d'Atiyah pour $\Omega^{1}$ \ill\ ; Koszul, 2 articles ;
\struck{Borel : lettre, notes}

\textbf{Deligne} \emph{(suite)} — Monsky--Washnitzer ; \struck{Serre :
cohomologie des vecteurs de Witt et von Dieudonné en caractéristique $p$} ;
Preprint Griffiths : « On monodromy \ldots » ; \struck{Lettre \ill\ chapitre de
familles} ; \struck{Work, Jacobiennes} \ill\ ; \struck{Papiers d'Artin (13)} ;
\struck{bibliothèque, \ill\ arithmétiques} ; \struck{Manin, variétés
rationnelles (2 articles)} ; \struck{\ill\ sur conjecture de Weil}

\page{4}

\textbf{\struck{\ill}} — Schlessinger, \emph{Functors of Artin rings} ;
\struck{Mes notes sur cat.\ de Picard}

\textbf{\struck{Saavedra}} — \struck{\ill\ de Chern des représentations de
groupes \ill} ; \struck{Note sur cat.\ de Picard} ; \struck{Rédaction \ill\
Picard enveloppant d'une catégorie triangulée} ; \struck{Notes catégories de
Picard} ; \struck{Notes \ill\ catégories} ; Note filtrations et groupes
paraboliques

\textbf{Mme Raynaud} \emph{1} — Papiers de Horrocks (Zariski) ;
Papier de \uncertain{Popp} (?) sur th.\ de Lefschetz etc.\ pour $\pi_{1}$ ;
Deligne, App.\ SGA 7 \struck{XXI}

\textbf{Mme Hakim} \emph{1} — mes notes sur tapis Quillen I, II

\textbf{\struck{Katz}} — \struck{notes sur connexions et équations de
Picard--Fuchs} ; \struck{notes sur torsion dans cohomologie cristalline, et sol.\
\ill\ de fond \ill\ sur cohomologie évanescente} ; \struck{notes EGA V}

\page{5}

\textbf{Deligne} \emph{2} \add{le chiffre est suivi d'un mot raturé, illisible} — \struck{Compléments \uncertain{Kodaira}} ; \struck{Mumford,
\emph{Amer.\ J.\ Math.} 1964} ; \struck{Mumford, \emph{Abelian quotients of
Teichmüller group}} ; Mumford, \emph{Some aspects of the Pb.\ of Moduli}
(miméo.) ; \struck{Lang--Néron} ; Matsusaka : \emph{Chow of Jac.\ variety},
\emph{A correction to\ldots}, \emph{The criteria for alg.\ equivalence} \ill\ ;
\struck{Néron : \ill\ fonctions et hauteurs sur les V.A.} ;
\struck{Kleiman : notes pour SGA 6} ; \struck{Grothendieck : lettre à Hartshorne
du 21.9.67 (sur ses papiers de dim.\ cohomologique des variétés algébriques)} ;
\struck{Weil} ; Andreotti \add{sur} Torelli ; \struck{Rédaction thèse \ill\ de
Abhyankar}

\textbf{\struck{Artin}} — \struck{Kiehl, sur th.\ de Frisch, \emph{Inventiones
Mathematicae}}

\textbf{Samuel} — Papiers de \uncertain{Ramanan}

\textbf{Berthelot} \add{(2)} \emph{1} — Gabriel : groupes formels ; Cartier :
notes \ill\ groupes formels ; notes sur stratifications ; \struck{notes sur
groupes de \uncertain{Barsotti--Tate} et cristaux} ; \struck{Kiehl, \ill}

\textbf{\struck{Artin}} / \textbf{Saavedra} — \struck{Atiyah, \emph{Equivariant
K-theory for compact Lie groups}} ; \struck{SGA 4 IX \add{?}} \ill

\textbf{\struck{B. Mazur}} — \struck{Correspondance avec Griffiths} ;
\struck{Notes Kleiman sur \uncertain{standard}}

\note{les « notes Kleiman sur standard » sont vraisemblablement celles sur les
conjectures standard ; le registre n'en dit pas davantage et on n'en dit pas
davantage ici}

\textbf{Deligne} \emph{(suite)} — \struck{Notes \ill\ avec Piatetski-Shapiro} ;
\struck{Lettre : Artin \ill\ cohomologie \ill}

\textbf{\struck{Messing}} — \struck{Conjecture sur \uncertain{Rosenlicht} \ill} ;
\struck{SGA 7 I \ldots\ V}

\textbf{\struck{Rim}} / \textbf{Quillen} — \struck{Rim, Gabriel} ; lettre à
Verdier ; Notes Verdier, limites projectives ; Notes Karoubi sur définition de
$K$.

\page{6}

\textbf{Szpiro} — Notes de moi et \uncertain{Bou.} sur caractérisation des
Gorenstein ; Notes sur $\exists$ complexe dualisant en dimension 1 ; Notes
SGA 7 ; \struck{Notes sur dualité de Poincaré en cohomologie de Hodge des
variétés algébriques singulières ((A) à (E))}

\textbf{\struck{Hartshorne}} — \struck{SGA 2} ; \struck{Livre Verdier}

\textbf{Balasko} — \struck{Notes : $NS_{X/S} \to NS_{B/S}$ etc., produit de
convolution sur V.A.} ; \struck{3 articles de Rim--Buchsbaum} ; \struck{Notes
EGA V sur idéaux de Fitting}

\textbf{\struck{Giraud}} — \struck{Notes Cartier sur vecteurs de Witt}

\textbf{\struck{\uncertain{Theodorou}}} — \struck{Bourbaki : \emph{Top.\ gén.}~I,
\emph{Alg.}~I}

\textbf{Poenaru} — Atiyah--Hirzebruch, \emph{über Cohomologie-Operationen} \ill\
; Stolzenberg : revue \add{du} livre Bishop \ill

\textbf{\struck{Boutot}} \emph{(suite)} — \struck{2 articles d'Artin sur
singularités rationnelles}

\textbf{\struck{Horrocks}} — \struck{Artin, \ill\ foncteurs représentables}

\textbf{Saint Donat} \emph{1} — SGA 5 IV \add{(}notes Jouanolou\add{)} ;
\struck{SGA 7 I \ldots\ V} ; Livre Mumford V.A. ; Notes Deligne \add{avec, en
regard :} th.\ de Lefschetz vache sur anneau pour courbes elliptiques
(cf.\ SGA 4 \uncertain{XVIII})

\page{7}

\textbf{Berthelot} \emph{2} \marginal{redemander}

— Deligne :
\struck{Structure : l'$\infty$ des \ill\ P.F.} ; \struck{lettre à Tate et
\uncertain{Washington}} ; Katz : th.\ de Dwork (2 papiers), \emph{Theory of
regular singular points}, \emph{Nilpotent connections} \add{et} Gauss--Manin
(2 \uncertain{versions}), th.\ de monodromie \add{remis à} Berthelot ;
Grothendieck : \emph{Torsion en cohomologie cristalline} ; Petit papier Deligne
sur prémorphismes de sites ; Papiers Deligne : Sém.\ Harvard ; SGA 5
\struck{XII} §§ 9, 10 \add{(}\ill\add{)} ; Notes Katz, connexions intégrables
(MS)

\textbf{Illusie} \textbf{\struck{\ill}} — lettre à Giraud sur obstructions aux
relèvements de groupes \uncertain{discrets} ; Papier Mazur--Roberts sur
cohomologie d'E.-P. ; \struck{2 paquets notes sur opérations de Cartier}
\add{« redemander », entouré} ; Notes sur th.\ de Deligne sur « diviseurs
universels » ; \struck{SGA \struck{4} 6 I, III, première rédaction}
\add{(car.\ 0)} ; Serre : variétés qui ne se relèvent pas en \ill

\textbf{\uncertain{Seydi}} — Notes sur lemme Hironaka ; \struck{Thèse Verdier}
sur morphismes universellement ouverts

\textbf{Thomas} \emph{1} — Papiers Lindenstrauss, Pełczyński etc.\ ; Survey de
Lindenstrauss \ill

\textbf{\uncertain{Zuckerman}} — Knorr, \emph{Finitude} I, II ; Petit livre sur
finitude

\textbf{Goldschmidt} — Brieskorn : singularités rationnelles \ill

\textbf{\struck{\ill}} — \struck{Artin, \emph{Algebraization of formal moduli}
I, II}

\page{8}

\textbf{\struck{\ill}} — \struck{2 lettres Deligne sur connexions intégrables
(sur $\mathbb{C}$)} ; \struck{Notes Serre, formes quadratiques} ;
\struck{Rapport Griffiths sur périodes} ; \struck{Séminaire Chevalley, groupes
par algèbres} ; \struck{Schlessinger, thèse} ; \struck{Tous tirages à part de
Bott} ; \struck{Borel : Springer Lecture Notes} ; \struck{Rim : SGA 7 VI}

\textbf{Demazure} — \struck{notes sur Kostant, sur résolution d'un morphisme
\ill}

\textbf{Saavedra} — \struck{Griffiths, \uncertain{Columbia}} ; \struck{Bourbaki,
Groupes de Lie} ; \struck{Notes sur polarisations dans catégories tannakiennes} ;
Notes sur structure de Hodge, lettre : Mumford, description transcendante des
motifs en caractéristique 0 ; \struck{Tate, cohomologie \add{des} tores} ;
\struck{Bourbaki, Algèbre}

\textbf{\uncertain{Bitzakis}} — \ill

\textbf{\struck{Van de Ven}} — Mes notes sur SGA 7 XX (Katz : le th.\ de
Griffiths) ; \struck{Griffiths : \emph{Report on Variation of Hodge
structure}}

\textbf{Saint Donat} \emph{2} — SGA 5 \struck{XVIII} ; Notes Mumford sur
variétés définies par équations quadratiques ; Livre Mumford V.A.

\textbf{Hubbard} — tirage à part sur rigidité des structures complexes
(\emph{Man.\ Math.})

\textbf{\uncertain{Nakis}} — Produits tensoriels topologiques (Grothendieck)

\page{9}

\textbf{Lascoux} — Borel--Hirzebruch I, II ; Borel, thèse ; Serre, Algèbre

\textbf{Mazur} — Lettre à Serre, corps de classes local etc.\ ; \struck{Notes
Katz, connexions « intégrables »} ; Quillen : \emph{Complex cobordism $+$ formal
groups} ; lettre de Bass \ill\ ; Rim, SGA 7 VII

\textbf{Teissier} — 2 articles Brieskorn, \ill\ d'un morphisme

\textbf{Figueroa} — Article Serre 1957

\textbf{\uncertain{Fouzat}} — Article Mumford, cycles

\textbf{Saavedra} — Groupes \ill, champs de vecteurs symétriques, \add{et}
Algèbres de Lie \add{Bourbaki} ; SGA 4 IV

\end{document}
